\section{Discussion and Future Work}
\label{sec:discussion}

\paragraph{What a static graph misses.}
H1--H3 together say that a single interaction graph leaves a measurable, predictable share of the
signal unused. Interactions depend on the program (H1, and the variance decomposition) and on how
far optimization has progressed (H2). Both effects can be learned from other programs (H3). The
drift is concentrated in the first simplification stage. This suggests that graphs should at
least be measured on simplified IR rather than on raw front-end output.

\paragraph{Why the method did not follow.}
H4 failed. The exploratory analysis points to a specific cause: graph-to-sequence builders score a
whole sequence by summing pairwise estimates. That ranks only a narrow family of near-identical
sequences, so additional budget is wasted. Better information (re-prediction, or even the oracle
graph) cannot fix a builder that does not explore. Search methods explore well but start from
nothing: at $\E=1$ they are far behind predicted graphs.

\paragraph{Proposal: prediction-seeded search.}
These observations suggest combining the two: a search, such as a GA or OpenTuner, whose initial
population is seeded with the top sequences of the program's predicted graph, and whose mutations
are drawn from that graph. Such a search would start where predict-once is strong and keep
improving the way search does. A static-graph variant of the same search would isolate the value
of per-program prediction. We have \emph{not} evaluated this design; it is a proposal. Testing it
would need its own pre-registration and held-out programs, because the idea came from the results
reported here.

\paragraph{Size versus speed.}
Optimizing IR instruction count is a code-size objective. The GA's sequences at $\E=200$ have
\MGAOThreeSizePct\% smaller x86-64 \texttt{.text} than \Othree{}. However, they execute
\MGAOThreeDynWorse\% more instructions and run \MGAOThreeRuntimeWorse\% slower. They have
\MGAOzIRPct\% fewer IR instructions than \Oz{}, yet \MGAOzSizeWorse\% larger x86-64 \texttt{.text}.
We attribute part of that gap to backend size heuristics: \Oz{} IR carries the \texttt{minsize}
attribute, which the backend honours. Our results therefore speak to IR- and code-size-oriented
ordering, the setting of embedded targets, not to speed.
